\section{开放挑战与瓶颈}
世界模型代表物理人工智能的前沿，但通往真正自主性的道路，受到感知合成与物理执行之间系统性裂缝的根本制约。当前范式能在互联网规模的被动视觉数据上轻松扩展，却在真实世界的交互性、安全关键约束面前十分脆弱。本节从整体视角分析世界模型面临的开放挑战：从使模型“视觉丰富却物理盲目”的结构性数据不对称（data asymmetry），到表面保真与可行动精度的背离，再到开环推演不断累积的漂移。我们认为，克服这些瓶颈需要根本性的范式转变：超越将世界模型仅作为生成渲染器的做法，迈向具有物理依据的模拟器和闭环交互规划器。这不仅需要考虑物理的归纳偏置与持续的环境数据采集，还需要重新构想评估框架、三方安全保证，以及能够协调集体智能与本地物理自主权的去中心化治理模式。

\subsection{数据不对称}
在渲染器、模拟器和规划器的功能分类下\citep{worldlabs2026taxonomy}，世界模型面临明确的数据不对称。渲染器可以依靠互联网规模图像、视频数据扩展，推动图像和视频生成快速进步\citep{kong2024hunyuanvideo,gao2025seedance,seedance2026seedance2,happyhorse2026}。但这些数据大多是被动、视觉性的：记录世界看起来怎样，却很少提供行动、物理属性、接触动力学或任务结果方面的监督。

相反，模拟器与规划器依赖结构更明确、可用于行动的数据。对模拟器而言，瓶颈在于可直接用于模拟的三维资产。虽然三维生成\citep{zhao2025hunyuan3d,hunyuan3d2025hunyuan3d,li2025triposg}迅速增加了视觉资产数量，但生成物体大多只适于渲染，而不能直接参与物理运行，往往缺少尺度、材料、碰撞几何、质量、摩擦、关节结构与可供性。PhysX-Anything\citep{cao2026physx}、PhysX-Omni\citep{cao2026phyomni}等近期研究，以及Lightwheel\citep{lightwheel2026}和Genesis World\citep{genesisworld2026}等工业平台，直接针对这一缺口；但真实感、非刚体动力学、长尾交互和可扩展性仍是开放挑战。对规划器而言，瓶颈进一步转向行动条件交互数据，如机器人演示、状态—行动轨迹和任务结果。Cosmos 3等近期模型\citep{nvidia2026cosmos3}将行动序列显式引入全模态世界建模，凸显行动条件数据对物理人工智能的重要性；但这类数据的可扩展性与透明程度，仍远逊于互联网视频。合成数据与模拟\citep{chang2017matterport3d,savva2019habitat,puig2024habitat,dosovitskiy2017carla,gan2020threedworld,gao2026nvidia}可部分缓解瓶颈，却引入视觉、几何、物理、行为和任务层面的分布变化。这些限制共同形成具身人工智能数据金字塔：底部是大量被动视频，中部是数量较少、可用于模拟的资产与交互轨迹，顶部则是稀缺却价值很高的触觉和力反馈。

应对这些系统性限制，需要根本改变物理智能获取世界信息的方式。现有数据采集受到规模边界与分布变化的内在限制，因此，未来世界模型的数据采集架构必须转向普遍存在、非侵入、持续运行的框架。首先，数据采集应无处不在，在不同物理场景中利用广泛、异构的采集机制，充分捕捉长尾环境。其次，数据获取流程应完全非侵入，在后台无缝运行，使采集不成为参与者的负担，也不打断日常工作；理想情况下，这种环境式采集还应通过人机协作主动提高运行效率。最后，采集必须持续进行，建立不间断的真实交互反馈循环，实时、动态、迭代地更新底层模型。

\subsection{保真度与精度}
世界模型的一项关键挑战，是感知保真度与物理精度不匹配。近期评测与研究给出了这类不匹配的经验结果。物理推理评估\citep{kang2024far,motamed2025generative,zhang2025morpheus,meng2024towards,gu2025phyworldbench,zheng2025vbench}表明，即使视觉质量提高，视频生成模型仍难以遵守物理规律和守恒约束，也难以实现分布外泛化。具身评测\citep{jiang2026robowm,shang2026worldarena}进一步揭示，视觉逼真的预测不一定转化为物理合理或可执行的机器人行为，暴露出用于操作与规划的世界模型在感知与功能之间的差距。

近期工作\citep{zhao2026phyworld,chen2026abot,wang2025physcorr,lin2026mmphysvideo}利用物理奖励、偏好配对、直接偏好优化（direct preference optimization，DPO）式训练，或结构化物理条件，在保留视觉质量的同时改善物理合理性。但当前奖励和偏好信号，仍是真实物理正确性的不完整替代，尤其难以涵盖细粒度接触、空间一致性和长尾交互。因此，感知保真度与物理精度之间的差距，仍是一项重要开放挑战。

要应对这一不匹配，需要从保真度优先的生成，转向以物理精度为中心的世界建模。首先，预训练应引入可扩展、物理信息丰富的监督，将视觉结果与潜在物理状态连接起来，如三维几何、物体运动、接触、力、材料属性和任务结果。其次，后训练应超越视觉偏好训练，引入度量约束满足、反事实有效性、长时域持续一致性和任务成功的物理奖励与评估。最后，应通过真实或高可信物理环境的闭环反馈校准世界模型，将预测与实际结果比较，并迭代改进。目标不仅是生成逼真未来，还要使世界模型预测在物理干预下能够指导行动、可以检验且可靠。

\subsection{预测误差累积}
即使世界模型局部准确，长时域部署仍可能失败，因为微小单步误差会递归反馈至后续预测。自回归推演中，每个预测观测、潜在状态或信念状态，都成为下一步上下文的一部分。因此，几何、速度、接触时刻、物体恒存性或隐藏状态估计中的细小误差会累积，将想象轨迹推离数据流形，最终偏离物理可实现行为。这一现象早已被认为是基于模型的强化学习与规划中模型偏差的主要来源\citep{talvitie2017selfcorrecting,lambert2022compounding}。视频预测与自回归视频生成也充分记录了类似失效模式：暴露偏差与时间漂移损害长时域推演\citep{nair2019hvf,chen2024diffusionforcing,huang2025selfforcing}。这同样是决策问题，而不仅是预测问题：规划器一旦开始针对不完美模型优化，就可能利用模型误差，将在真实环境中绝不可能成功的轨迹赋予高价值。在Dyna类规划中，这被强调为“幻觉出价值”\citep{jafferjee2020hallucinating}。在视觉丰富、接触密集的领域，挑战尤其突出：部分可观测性、分支未来和动力学中的剧烈不连续，使长开环推演内在脆弱。

多种缓解策略已被证明有用，但没有一种消除了问题。第一条路线是在紧凑潜在状态空间而非原始像素中预测，减少无关感知噪声，迫使模型聚焦决策相关动力学，如PlaNet和Dreamer家族\citep{hafner2019planet,hafner2019dreamer,hafner2021dreamerv2,hafner2023dreamerv3}。第二条路线通过概率动力学模型、集成和保守推演时域，使规划考虑不确定性，让行动选择更少依赖单一、脆弱的想象未来，更多依赖对可能未来经过校准的不确定性\citep{chua2018pets,janner2019mbpo,hansen2022tdmpc,hansen2023tdmpc2}。第三条路线通过自纠正目标、层级抽象和多尺度规划改善时间一致性，周期性地让长时域推理回到观测依据，而非无限期信任单一开环推演\citep{talvitie2017selfcorrecting,hafner2022director,gumbsch2024learning}。这些方法共同表明，关键目标不只是降低单步预测损失，而是让世界模型在真实规划所需时域内，保持误差有界、经过校准且适合决策。

\subsection{仿真到现实迁移}
在机器人中部署世界模型的根本挑战，是所学模拟与真实物理之间的差距。世界模型虽旨在从数据近似环境动力学，但即使微小建模误差，也可能在长时域推演中快速累积，造成不断加剧的分布变化，使真实世界预测不可靠。机器人中的这一问题尤其严重，因为接触丰富的动力学、部分可观测性及高维行动空间，使准确长期预测极其困难。

为缓解仿真到现实差距，人们探索了多种策略。经典方法是域随机化（domain randomization）：训练时让策略或动力学模型接触多样化模拟变化，促使其对视觉与物理扰动具有不变性，提高鲁棒性\citep{tobin2017domain,peng2018sim}。另一条路线关注系统辨识（system identification），从有限真实观测推断未知物理参数，用于校准模拟器或所学动力学模型\citep{nagabandi2018learning}。近期的在线适应方法，使世界模型能够利用真实反馈持续更新动力学，从而减少分布变化下的预测误差\citep{hafner2019dreamer,nagabandi2018neural}。

对基于视频的世界模型和世界行动模型（WAM）而言，这一挑战更突出：它们在高维像素空间运行，高度依赖大规模预训练所学视觉动力学先验。虽然这些模型受益于对多样互联网视频的强泛化能力，但由于物理规律、机器人本体和行动含义存在差异，将这些先验迁移至精确机器人控制仍非易事。此外，仅依靠视觉观测往往不足以支持机器人，因为视频缺少底层物理状态、本体感觉反馈及细粒度行动动力学的显式信息。加入本体感觉、触觉、力反馈及结构化机器人状态等额外模态，可为环境动力学提供互补约束，帮助缩小仿真到现实差距。应对这一差距，需要更好地整合考虑物理的归纳偏置、结构化状态表征，以及推理时自适应修正机制；这些仍是世界模型在真实机器人中规模化部署的开放挑战。

\subsection{评估与基准测试}
世界模型的评估格局仍很分散。世界模型跨越视觉预测、物理模拟、行动规划和下游决策，不同研究群体常使用不同指标与评测。对预测及生成模型，常用指标包括均方误差（mean squared error，MSE）等配对重建误差，以及FID、FVD等分布或感知指标，度量生成图像、视频相对于真实数据的质量。对控制与规划任务，通常依据任务表现评估，如奖励、回合回报、成功率或任务完成情况。泛化通常在分布外（OOD）或零样本条件下评估，在未见环境、物体、任务或动力学上度量相同预测或任务指标。在机器人和自动驾驶等安全关键场景，不确定性校准等面向可靠性的指标同样重要，因为智能体必须知道何时预测或计划不可靠，才能安全推演轨迹并进行考虑风险的规划。

在评测层面，这些指标用于一系列代表性基准测试。对预测式与生成式视频模型，VBench\citep{huang2023vbench}和VBench-2.0\citep{zheng2025vbench}等套件评估视频质量、时间一致性、与提示的符合程度及内在忠实性；与Cosmos 3一同引入的Cosmos-HUE等人工评估规程\citep{nvidia2025cosmos}，评估生成式与全模态世界模型的输出。对三维／四维世界生成，WorldScore\citep{duan2025worldscore}和4DWorldBench\citep{lu20264dworldbench}将评估从二维视频质量扩展到可控性、动力学和时空一致性。对物理一致性及世界模拟器保真度，PhyGenBench\citep{meng2024towards}、WorldModelBench\citep{li2026worldmodelbench}和PhyWorldBench\citep{gu2025phyworldbench}等面向物理的评测，考察生成视频是否遵循物理常识、具有物理真实感；WorldSimBench\citep{qin2024worldsimbench}进一步从感知和行动两个层面，将视频生成模型作为世界模拟器评估。

对行动条件与下游功能评估，Atari／ALE\citep{bellemare2013arcade}和DM Control\citep{tassa2018deepmind}常用于基于模型的强化学习和连续控制；Habitat\citep{savva2019habitat}和CARLA\citep{dosovitskiy2017carla}在模拟中评估具身导航和自动驾驶；RoboArena\citep{atreya2025roboarena}将评估扩展到真实机器人策略；EWMBench\citep{yue2025ewmbench}从场景一致性、运动正确性与语义符合程度评估具身世界模型；WorldArena\citep{shang2026worldarena}评估其功能效用。对于交互式世界模型，WBench\citep{ying2026wbench}和WorldMark\citep{xu2026worldmark}等近期评测，考察模型能否遵循多轮交互、保持控制一致性和物理合理性。作为补充，MBench\citep{zhang2026mbench}、WorldPrediction\citep{chen2025worldprediction}和WorldReasonBench\citep{wu2026worldreasonbench}等长时域状态推理评测，考察世界模型能否维持持续状态、支持程序性规划，并生成在物理、社会、逻辑与信息层面具有一致状态演化的未来视频。

近期CoW-Bench进一步围绕“一致性三位一体”——模态、空间与时间一致性——组织世界模型评估，测试生成模型在复杂多模态情境下，能否保留跨模态对应关系、具有三维感知的空间结构、物体恒存性，以及物理合理的时间演化\citep{wei2026trinityconsistencydefiningprinciple}。尤其需要指出，现有评测往往忽略推理效率；但真实物理交互所需的闭环控制频率，无法容忍高延迟推演。亚秒级生成速度不只是计算指标，也是基本安全要求，因为推理滞后会直接引发高动态环境中的控制不稳定。

归根结底，当前评估范式的根本瓶颈，是缺少真实物理依据。虚拟模拟器采用简化物理引擎，使世界模型能够利用在现实中失效的非因果捷径。真正的验证需要真实部署和持续物理交互，以面对未被建模的随机动力学。然而，转向物理平台也引入了真实性与实验控制程度之间的重要取舍。核心开放挑战是评测标准化与公平性。与可复现的软件容器不同，物理评估受到硬件退化、初始状态不一致，以及极高重复成本的影响。如何制定规程，规范不同机器人本体的评估，并将硬件噪声与模型能力分离，仍是下一步的重要研究方向。

\subsection{安全、透明性与可持续性}
\paragraph{安全探索与三方交互安全。}
具身人工智能和自主智能体的部署引入显著物理风险，因为在不受约束的环境中实时决策，比孤立的数据分析危险得多。模拟与现实的残余差距永远无法彻底消除，因此，未经测试就在真实环境探索，本身便具有危险\citep{billard2025roadmap}。物理世界模型充当学习得到的内部模拟器，使智能体策略能在部署前通过想象推演充分优化，从而缓解风险。推理时，这些模型提供连续状态表征与多步规划，以预判行动后果。与重型或高吞吐量机械交互时，为保证安全，必须在推演引擎中嵌入约束方法和严格的不确定性阈值；一旦预测外推至不安全状态区域，就停止执行。由于环境、机器人本体和任务会在智能体运行生命周期中演化，安全探索必须是一项与终身学习耦合的持续能力。

此外，这种安全范式会在动态三方生态中非线性扩展，涉及人机协调、环境反馈的共同演化，以及异构多智能体同步。首先在人机交互中，安全不止于反应式避障，而成为博弈论意义上的认知协调；人类应对行为会根据机器人轨迹动态变化，因此世界模型必须将心智理论框架融入潜在推演，主动适应隐含人类意图。其次，在机器—环境循环中，世界模型必须体现因果双向性：不仅预测智能体如何改变周围环境，还要预测非平稳环境反馈，如结构疲劳或热变形，如何反过来破坏控制保证。第三，不同形态的异构智能体共处时，主要风险会从单个控制器失效，转向系统性共振：各自世界模型之间不对称的不一致，可能级联成分布式碰撞。

\paragraph{透明性、可解释性与可验证控制。}
要合乎伦理、可靠地部署世界模型，可解释性对建立操作人员信任和厘清责任至关重要。与传统黑箱端到端策略相比，世界模型依赖环境动力学的显式表征，因而提供了结构更明确的透明化路径。但深度神经网络本身难以解释，将这些动力学转化为可验证预测，仍是开放工程问题。仅依赖人为构建的可解释层，维持含义清晰的潜在维度，以运行时验证速度界限、能量守恒等不变量，还不够。所学组件必须系统地结合控制理论或形式化验证框架，取得可证明的稳定性与约束满足，使预测具有数学保证，而非不受约束的外推。

此外，操作人员进入循环后，单靠技术可解释性无法消除系统风险。世界模型高度逼真的模拟输出，可能引发自动化偏差（automation bias）等认知弱点；已有研究记录，操作人员容易过度依赖看似权威的模型预测，即使这些预测错误或处在分布外\citep{parasuraman1997humans}。这一弱点源于人如何在认知上接受解释，而非底层模型准确性，因此仅改善可解释性无法解决。缓解自动化偏差，需要面向用户的界面显式呈现认识不确定性，例如标记偏离已验证训练分布的推演，使人的信任准确对应系统的客观可靠程度。

\paragraph{合乎伦理的部署与保护隐私的治理。}
合乎伦理的部署是一项独立要求，涉及根据世界模型的人类相关预测与数据采取行动是否合法正当。这带来三项中心义务。第一，预测人类行为以规划机器人行动时，世界模型不得依赖基于偏见的画像\citep{billard2025roadmap}。这一约束极难落实，因为预测直接进入物理行动，而非供人拒绝的信息建议。第二，数据采集存在严格隐私边界：物理交互数据集中的参与者必须不可识别，不应能推断敏感属性，记录也不得超出原始范围转作他用\citep{billard2025roadmap}。这些约束虽必要，却严重限制大规模人类交互数据集的构建。第三，透明性是事故责任认定的前提，但商业上普遍采用闭源感知和控制系统，抵制独立审查，这与透明要求直接冲突\citep{billard2025roadmap}，使问责仍是开放治理挑战。

为解决跨机构数据汇集障碍并兼顾商业保密，协作框架必须转向联邦世界模型（federated world model）。在这一范式下，本地数据孤岛保持原状，参与方只交换潜在梯度或网络权重。为防御梯度反演攻击（gradient inversion attack），这些架构必须整合安全多方计算（secure multi-party computation），或使用硬件强制保护的可信执行环境（trusted execution environment，TEE）进行加密联合训练。关键在于，为适应不同机构环境中的非独立同分布（Non-IID）数据\footnote{译注：原文将Non-IID展开为“non-identical and independently distributed”，与该缩写通常表示的“非独立同分布”不符；此处按Non-IID及上下文翻译。}，框架必须将不变的物理常识——全局学习、普遍共享，以支撑集体智能——与变化很大的领域特定潜在表征——本地保留——分离，从而兼顾集体智能与本地数据主权。

不过，还有一项重要限制：当前优化几乎只关注训练和推理运行阶段的能耗，完全忽略硬件整个生命周期的环境足迹。完整的可持续性评估必须纳入生命周期管理框架，涵盖材料开采、制造与报废处理\citep{billard2025designing}。稀土采矿、软执行器聚合物合成和半导体制造等上游环境成本，在当今人工智能评测中大多被忽略。可持续计算文献表明，隐含碳排放（embodied carbon）与运行碳排放（operational carbon）往往存在取舍；孤立优化其中一项，可能无意中放大另一项\citep{lee2025view}。因此，仅依据计算指标提出效率结论，从根本上说是不完整的，必须审慎解读。
